\section{The Integer Base}
\label{sec:base}

The geometry of Section~\ref{sec:geometry} is not chosen. It is forced from the least arbitrary
starting point there is, the whole numbers. This section walks the ladder that climbs from the
integers to the crystal (P47 to P51).

\paragraph{The modular group (P48).}
The symmetries of the integers, the structure-preserving rearrangements, form the modular group,
generated by two simple moves and free of any tuned parameter. The testbed builds it directly from
the integers, with no constants put in. This is the algebraic seed of everything above it.

\paragraph{The Coxeter construction (P47).}
Widening symmetry to reflection turns the modular structure into a Coxeter group, and a Coxeter
group acting on hyperbolic space produces a regular tessellation. The only input is the integer
angles between mirrors, the Schl\"afli symbol. So the crystal of Section~\ref{sec:geometry} is the
geometric face of an algebraic object that is itself the symmetry of the integers.

\paragraph{The golden thread (P50).}
The golden ratio $\varphi = (1+\sqrt 5)/2$ appears in the foundation from independent directions:
it is the growth ratio of the crystal ring by ring, the most evenly spread scatter of points, and
the straightest path through the Fibonacci addressing. Three separate roads arrive at the same
constant, a sign of a single ordered structure rather than a patchwork. Its integer shadow, the
Fibonacci sequence, is the addressing and ring-counting scheme of P42.

\paragraph{The full ladder (P51).}
P51 runs the whole climb in one program: integers, then their symmetry group, then the reflection
group, then the hyperbolic crystal, then the tilings and honeycombs, with nothing arbitrary added
at any rung. The companion note on the integer ladder gives the wider-audience account. The upshot
for this paper is that the substrate's geometry is a theorem about the integers, not a modelling
choice, and the dimension classification of Section~\ref{sec:geometry} is the same statement read
geometrically.

\paragraph{What is special among the integers.}
A separate analysis records which integers are load-bearing. The small ones carry the laws, because
a law needs the smallest sufficient structure: three is the least alphabet with a negative, a
neutral, and a positive, hence the ternary tone. Seven is the least $p$ that curves space at
$q = 3$, hence the heptagrid. The dimension window is two-to-four. Large integers are only
addresses, all equal. The framework is not prime-centric. Importance is set by role, not by
magnitude.
